\section{Discussion}
\label{sec:discussion}

\paragraph{The bottleneck is the revision, not the gate.} The analytical model separates two questions that are easy to confuse: whether revising helps on the questions it touches, measured by the fix and break rates, and whether a confidence signal can find those questions, measured by its ROC curve. Our results show that the first question decides the outcome. With self-critique, the fix rate is zero on GSM8K and very small on HotpotQA, so no signal can make revision useful, even though self-consistency identifies wrong answers well. With majority voting as the revision step, the same signals and the same gates produce gains that match the model's prediction. A gate can remove harm, but it cannot create fixes.

\paragraph{Why self-critique fails here.} The revised answers show a clear pattern of self-confirmation: the model repeats its earlier reasoning and declares it correct. Part of this is probably caused by our revision prompt, which explicitly allows the model to keep its answer when it finds no error. A more critical prompt might raise the fix rate, but, as the HotpotQA results suggest, it would probably raise the break rate as well. Measuring both rates is more informative than reporting net accuracy alone.

\paragraph{Answers to the research questions.} \textbf{RQ1:} With self-critique as the revision step, no gate outperforms never revising; with majority voting as the revision step, gating keeps most of the benefit of majority voting while revising far fewer questions, but its advantage over never revising is not significant on our test sets, and a small development split can fail to find a useful threshold. \textbf{RQ2:} Self-consistency is the most discriminative signal (AUROC 0.95 on GSM8K and 0.69 on HotpotQA); verbalized confidence and P(True) are weak (0.58--0.64). \textbf{RQ3:} The best signal is also the most expensive: self-consistency needs about eleven times the tokens of a single answer, whereas P(True) is cheap to generate but expensive to prefill when the context is long.

\paragraph{Comparison with prior work.} Our always-revise results agree with Huang et al.~\cite{huang2024cannot}: without external feedback, revision does not improve reasoning accuracy and can reduce it. In the terms of Yang et al.~\cite{yang2024confidence}, \modelA{} shows high confidence (it rarely changes a correct answer) but very low critique ability (it almost never repairs a wrong one). IoE~\cite{li2024confidence} behaves like a cautious revision step: it changes few answers and, on HotpotQA, still breaks more than it fixes. The strength of majority voting is consistent with Wang et al.~\cite{wang2023selfconsistency}.

\paragraph{Practical recommendation.} Before adding a confidence gate to a self-correction pipeline, practitioners should estimate the fix and break rates of the revision step on a development set. If $(1-a)f$ is small, the expected benefit of any gate is small as well, and effort is better spent on a revision step that actually repairs errors, such as resampling, or on obtaining external feedback.
